\documentclass[10pt,a4paper]{amsart}
\usepackage[margin=19mm]{geometry}
\usepackage{amsmath,amssymb,mathtools}
\usepackage[scr=rsfs,cal=euler]{mathalfa}
\usepackage{xfrac}
\usepackage[unicode,hidelinks,bookmarks=false]{hyperref}
\newcommand{\ie}{i.e.\ }
\newcommand{\cf}{cf.\ }
\newcommand{\resp}{respectively\ }
\newcommand{\Subsection}{\subsection}
\providecommand{\nobreakdash}{\nobreak\hbox{-}}
\setlength{\emergencystretch}{2em}
\multlinegap0pt
\usepackage{amssymb}
\usepackage{enumerate}
\usepackage{mathtools}
\newtheorem{lemma}{Lemma}
\newcommand{\bbs}{\mathbb S}
\newcommand{\calA}{\mathcal{A}}
\newcommand{\calF}{\mathcal{F}}
\newcommand{\dm}{d} % dimension
\newcommand{\dx}{\mathrm{d}S (x)}
\newcommand{\dz}{\mathrm{d}S (z)}
\newcommand{\supp}{\mathrm{supp}(\rho)}
\title{Stability of equilibria of an aggregation-diffusion energy on sphere}
\author{Razvan C. Fetecau, Hansol Park}
\date{Source: arXiv:2608.26415v1 (CC BY 4.0)}
\begin{document}
\maketitle

\noindent\emph{Source and licence.} Stability of equilibria of an aggregation-diffusion energy on sphere. Version arXiv:2608.26415v1 (CC BY 4.0), distributed by arXiv under the Creative Commons Attribution 4.0 International licence, http://creativecommons.org/licenses/by/4.0/, as recorded in the arXiv licence metadata for this exact version and shown on the abstract page https://arxiv.org/abs/2608.26415v1. Full licence notice: \texttt{licenses/arxiv-2608.26415-CC-BY-4.0.txt}.\par\smallskip

\noindent\textit{Editorial scope.} Counted unit: the article's lemma lem:approx (source0.tex lines 1559--1563). The article's complete original proof of it is delivered with it (source0.tex lines 1564--1738, one \texttt{proof} environment of 175 lines). No mathematics is written, repaired or re-derived: the statement and the proof are the article's own lines, in the article's own notation, and no retained line is substituted or re-flowed.  Retained uncounted as the source-local context the proof needs: lines 940--943; lines 1135--1137.  Nothing was omitted from any retained span, so the package carries no omission record.  No retained line contains a citation, so the wrapper carries no bibliography and no bibliography entry.  No retained line contains a figure environment, an image-inclusion command or drawing code, so no image is delivered and the package declares no asset dependency.  Editorial formatting only, in addition: an AMS \texttt{amsart} preamble that restates the article's own declarations and macros used by the retained lines, namely the environments \texttt{lemma} on separate counters, together with the standard packages those lines need.  The retained environments print in this document as Lemma 1 in order of appearance, whereas the article prints the same items with the numbering of its own full document.  The complete native source of the article as distributed in the arXiv e-print for this exact version is bundled unchanged as \texttt{sources/208632/source0.tex}.
\begin{equation}
\label{eqn:calA}
\mathcal{A}:=\left\{\psi \in L^1:\mathrm{supp}(\psi)\subset \supp,\quad\int_{\supp}\psi(x)\dx=0,\quad \psi(x)\geq-\frac{1}{\epsilon_0}\rho(x)\right\}.
\end{equation}

\begin{align}\label{eqn:phicy}
\psi_y(x)=\frac{\rho(x)^{2-m}y\cdot(x-c_{\rho^{2-m}})}{\int_{\supp}\rho(z)^{2-m}\left(y\cdot (z-c_{\rho^{2-m}})\right)^2\dz}.
\end{align}

\begin{lemma}\label{lem:approx} Let $y\in \bbs^\dm$ be either parallel or perpendicular to $x_0$. Then, there exists a sequence $\{\psi_{n, y}\}_{n\geq1}\subset \mathcal{A}$ such that
\[
\lim_{n\to\infty}\mathcal{F}[\psi_{n,y}]=\mathcal{F}[\psi_{y}].
\]
\end{lemma}

\begin{proof}
Let $y$ be a fixed vector either parallel or perpendicular to $x_0$. Recall \eqref{eqn:phicy}, and set
\[
\alpha:=\int_{\supp}\rho(x)^{2-m}\left(y\cdot (x-c_{\rho^{2-m}})\right)^2\dx.
\]
Define 
\[
\tilde{\psi}_y(x):=\alpha\psi_y(x)=\rho(x)^{2-m}y\cdot(x-c_{\rho^{2-m}}).
\]
Since the functional $\calF$ is homogeneous, we have
\begin{equation}
\label{eqn:Ftildephi=Fphi}
\mathcal{F}[\tilde{\psi}_y]=\mathcal{F}[\psi_y].
\end{equation}
As the function $\tilde{\psi}_y$ is much simpler in form and yields the same functional value as $\psi_y$, we will work with it instead.

For $R>0$ define
\[
\tilde{\psi}_y^R(x)=\rho(x)\min(\rho(x)^{1-m}, R)\left(y\cdot\left(x-c_{\rho\min(\rho^{1-m}, R)}\right)\right).
\]
By their construction, $\tilde{\psi}_y^R$ is expected to approximate $\tilde{\psi}_y$ for large $R$; this will be shown formally below. Also, $\tilde{\psi}_y^R$ satisfies
\begin{equation}
\label{eqn:int-tildephiy}
\int_{\supp}\tilde{\psi}_y^R(x)\dx=0,
\end{equation}
and furthermore,
\begin{equation}
\label{eqn:tildephiy-ineq}
\tilde{\psi}_y^R(x)\geq -2R\rho(x),
\end{equation}
as $\min(\rho(x)^{1-m},R)\leq R$ and 
\begin{align*}
\left|y\cdot\left(x-c_{\rho \min(\rho^{1-m}, R)}\right)\right| &\leq \|y\| \left( \|x\| + \| c_{\rho \min(\rho^{1-m}, R)}\| \right) \\[2pt]
&\leq 2.
\end{align*}
By \eqref{eqn:int-tildephiy}, \eqref{eqn:tildephiy-ineq}, and the definition of $\calA$ in \eqref{eqn:calA}, we have $\frac{1}{2R\epsilon_0}\tilde{\psi}_y^R\in\mathcal{A}$; this observation will be used later to define the sequence $\psi_{n,y} \in \calA$ needed to show the lemma.

To compare $\tilde{\psi}_y^R$ and $\tilde{\psi}_y$, we start by estimating the difference between $c_{\rho\min(\rho^{1-m},R)}$ and $c_{\rho^{2-m}}$. We compute
\begin{align*}
&c_{\rho\min(\rho^{1-m}, R)}-c_{\rho^{2-m}}\\[3pt]
&\quad = \frac{\int_{\supp}\rho(x)\min(\rho(x)^{1-m},R)x \, \dx}{\int_{\supp}\rho(x)\min(\rho(x)^{1-m},R)\, \dx}-\frac{\int_{\supp}\rho(x)^{2-m}x \, \dx}{\int_{\supp}\rho(x)^{2-m} \, \dx}\\[3pt]
&\quad = \frac{\int_{\supp} \rho(x)\min(\rho(x)^{1-m},R)x\dx \; \int_{\supp} \left(\rho(x)^{2-m}-\rho(x)\min(\rho(x)^{1-m},R) \right)\dx
}{\int_{\supp}\rho(x)\min(\rho(x)^{1-m},R)\dx \, \int_{\supp}\rho(x)^{2-m}\dx}\\[3pt]
&\qquad + \frac{\int_{\supp}\rho(x)\min(\rho(x)^{1-m},R)\dx \,\int_{\supp} \left(\rho(x)\min(\rho(x)^{1-m},R)-\rho(x)^{2-m} \right)x \dx
}{\int_{\supp}\rho(x)\min(\rho(x)^{1-m},R)\dx \, \int_{\supp}\rho(x)^{2-m}\dx}.
\end{align*}
From the triangle inequality, we then get
\begin{align*}
&\left \|c_{\rho\min(\rho^{1-m}, R)}-c_{\rho^{2-m}} \right\|\\[3pt]
&\leq\left\|\frac{\int_{\supp} \rho(x)\min(\rho(x)^{1-m},R)x\dx \, \int_{\supp} \left(\rho(x)^{2-m}-\rho(x)\min(\rho(x)^{1-m},R)\right)\dx
}{\int_{\supp}\rho(x)\min(\rho(x)^{1-m},R)\dx\, \int_{\supp}\rho(x)^{2-m}\dx}\right\|\\[3pt]
&\quad +\left\|\frac{\int\rho(x)\min(\rho(x)^{1-m},R)\dx\int_{\supp}\left(\rho(x)\min(\rho(x)^{1-m},R)-\rho(x)^{2-m} \right) x \dx
}{\int_{\supp}\rho(x)\min(\rho(x)^{1-m},R)\dx\int_{\supp}\rho(x)^{2-m}\dx}\right\|.
\end{align*}

If we use $\|x\|=1$ for all $x\in\bbs^\dm$, we can further estimate the r.h.s. above and find
\begin{align*}
&\left \|c_{\rho\min(\rho^{1-m}, R)}-c_{\rho^{2-m}} \right\|\\[3pt]
&\leq 2 \, \frac{\int_{\supp}\rho(x)\min(\rho(x)^{1-m},R)\dx \, \int_{\supp}|\rho(x)^{2-m}-\rho(x)\min(\rho(x)^{1-m},R)|\dx}{\int_{\supp}\rho(x)\min(\rho(x)^{1-m},R)\dx\int_{\supp}\rho(x)^{2-m}\dx}.
\end{align*}
By the dominated convergence theorem, we have
\begin{equation}
\label{eqn:2limR}
\begin{aligned}
&\lim_{R\to\infty}\int_{\supp}\rho(x)\min(\rho(x)^{1-m},R)\dx =\int_{\supp}\rho(x)^{2-m}\dx,\\
&\lim_{R\to\infty}\int_{\supp}|\rho(x)^{2-m}-\rho(x)\min(\rho(x)^{1-m},R)|\dx=0.
\end{aligned}
\end{equation}
Therefore, we can conclude that
\begin{equation}
\label{eqn:limR-normc}
\lim_{R\to\infty}\|c_{\rho\min(\rho^{1-m}, R)}-c_{\rho^{2-m}}\|=0.
\end{equation}

Now, we compare $\tilde{\psi}_y^R$ and $\tilde{\psi}_y$ directly:
\begin{align*}
&\tilde{\psi}_y(x)-\tilde{\psi}_y^R(x)\\[2pt]
&=\rho(x)^{2-m}y\cdot(x-c_{\rho^{2-m}})-\rho(x)\min(\rho(x)^{1-m}, R)\left(y\cdot\left(x-c_{\rho\min(\rho^{1-m}, R)}\right)\right)\\[2pt]
&=\rho(x)^{2-m}y\cdot(c_{\rho\min(\rho^{1-m},R)}-c_{\rho^{2-m}})-\rho(x)\left(\min(\rho(x)^{1-m}, R)-\rho^{1-m}\right)\left(y\cdot\left(x-c_{\rho\min(\rho^{1-m}, R)}\right)\right).
\end{align*}
Then, we estimate
\begin{align*}
&\|\tilde{\psi}_y-\tilde{\psi}_y^R\|_{1}\\
&\leq\int_{\supp}|\rho(x)^{2-m}y\cdot (c_{\rho\min(\rho^{1-m}, R)}-c_{\rho^{2-m}})|\dx\\
&+\int_{\supp}|(\rho^{2-m}(x)-\rho(x)\min(\rho(x)^{1-m},R))( y\cdot(x-c_{\rho\min(\rho^{1-m}, R)}))|\dx.
\end{align*}
We use again $\|y\|\leq1$ and $\|x-c_{\rho\min(\rho^{1-m},R)}\|\leq2$ to get
\begin{align*}
&\|\tilde{\psi}_y-\tilde{\psi}_y^R\|_{1}\\
&\leq \|c_{\rho\min(\rho^{1-m},R)}-c_{\rho^{2-m}}\|\int_{\supp}\rho(x)^{2-m}\dx\\
&+2\int_{\supp}|\rho^{2-m}(x)-\rho(x)\min(\rho(x)^{1-m},R)|\dx.
\end{align*}
By \eqref{eqn:2limR} and \eqref{eqn:limR-normc}, we finally find
\begin{equation}
\label{eqn:limphi-L1}
\lim_{R\to\infty}\|\tilde{\psi}_y-\tilde{\psi}_y^R\|_{1}=0.
\end{equation}

Next, we calculate the difference between $\mathcal{F}[\tilde{\psi}_y]$ and $\mathcal{F}[\psi_{y}^R]$:
\begin{align*}
&\mathcal{F}[\tilde{\psi}_y]-\mathcal{F}[\tilde{\psi}_y^R]\\[2pt]
&\quad =\frac{\int_{\supp}\rho(x)^{m-2}\tilde{\psi}_y(x)^2\dx}{\|\mu_{\tilde{\psi}_y}\|^2}-\frac{\int_{\supp}\rho(x)^{m-2}\tilde{\psi}_y^R(x)^2\dx}{\|\mu_{\tilde{\psi}_y^R}\|^2}\\
&\quad =\frac{1}{\|\mu_{\tilde{\psi}_y}\|^2}\int_{\supp}\rho(x)^{m-2}\left(\tilde{\psi}_y(x)^2-\tilde{\psi}_y^R(x)^2\right)\dx\\[2pt]
&\qquad +\frac{\|\mu_{\tilde{\psi}_{y}^R}\|^2-\|\mu_{\tilde{\psi}_y}\|^2}{\|\mu_{\tilde{\psi}_y}\|^2\|\mu_{\tilde{\psi}_{y}^R}\|^2}\int_{\supp}\rho(x)^{m-2}\tilde{\psi}_y^R(x)^2\dx.
\end{align*}
We further estimate
\begin{equation}
\label{eqn:F-diff}
\begin{aligned}
\left|\mathcal{F}[\tilde{\psi}_y]-\mathcal{F}[\tilde{\psi}_y^R]\right|&\leq
\frac{1}{\|\mu_{\tilde{\psi}_y}\|^2}\left|
\int_{\supp}\rho(x)^{m-2}(\tilde{\psi}_y(x)^2-\tilde{\psi}_y^R(x)^2)\dx\right|
\\
&+\frac{\left|\|\mu_{\tilde{\psi}_{y}^R}\|^2-\|\mu_{\tilde{\psi}_y}\|^2\right|}{\|\mu_{\tilde{\psi}_y}\|^2\|\mu_{\tilde{\psi}_{y}^R}\|^2}
\int_{\supp}\rho(x)^{m-2}\tilde{\psi}_y^R(x)^2\dx.
\end{aligned}
\end{equation}
The second term in the r.h.s. of \eqref{eqn:F-diff} tends to zero as $R\to\infty$ since
\[
\lim_{R\to\infty}\|\mu_{\tilde{\psi}_y^R}-\mu_{\tilde{\psi}_y}\|=\lim_{R\to\infty}\left\|\int_{\supp}(\tilde{\psi}_y^R(x)-\tilde{\psi}_y(x))x\dx\right\|\leq \lim_{R\to\infty} \|\tilde{\psi}_y^R-\tilde{\psi}_y\|_1=0,
\]
by using \eqref{eqn:limphi-L1}. We will also show that as $R \to \infty$, the first term in the r.h.s. of \eqref{eqn:F-diff} tends to zero as well. 

Indeed, as $x, y\in\bbs^\dm$ and $\|c_{\rho^{2-m}}\|,\|c_{\rho\min(\rho^{1-m},R)}\|\leq1$, we have
\[
|\tilde{\psi}_y(x)|,|\tilde{\psi}_y^R(x)|\leq 2\rho(x)^{2-m}, \qquad\forall x\in\supp,\quad R>0,
\]
from which we can infer
\begin{align*}
\rho(x)^{m-2}\left|\tilde{\psi}_y(x)^2-\tilde{\psi}_y^R(x)^2\right|&=\rho(x)^{m-2}\left|\tilde{\psi}_y(x)+\tilde{\psi}_y^R(x)\right|\times\left|\tilde{\psi}_y(x)-\tilde{\psi}_y^R(x)\right|\\
&\leq4\left|\tilde{\psi}_y(x)-\tilde{\psi}_y^R(x)\right|.
\end{align*}
Hence, we get
\begin{align*}
\left|
\int_{\supp}\rho(x)^{m-2}(\tilde{\psi}_y(x)^2-\tilde{\psi}_y^R(x)^2)\dx\right|& \leq 4\int_{\supp}|\tilde{\psi}_y(x)-\tilde{\psi}_y^R(x)|\dx \\
&=4\|\tilde{\psi}_y-\tilde{\psi}_y^R\|_1,
\end{align*}
and by \eqref{eqn:limphi-L1} we conclude that the first term in the r.h.s. of \eqref{eqn:F-diff} tends to zero as $R\to\infty$. Combining the results, we find from \eqref{eqn:F-diff} that
\[
\lim_{R\to\infty}\left|\mathcal{F}[\tilde{\psi}_y]-\mathcal{F}[\tilde{\psi}_y^R]\right|=0,
\]
and hence,
\begin{equation}
\label{eqn:limF-Rinf}
\lim_{R\to\infty}\mathcal{F}[\tilde{\psi}_y^R]=\mathcal{F}[\tilde{\psi}_y].
\end{equation}

Now, we choose a sequence 
\[
R_n\nearrow\infty,\quad\text{ as }\quad n\nearrow\infty.
\]

Since $\tilde{\psi}_y^{R_n}$ satisfies (see \eqref{eqn:tildephiy-ineq})
\[
\tilde{\psi}_y^{R_n}(x)\geq-2R_n\rho(x),\qquad\forall x\in\bbs^\dm,
\]
we define
\[
\psi_{n, y}=\frac{1}{2R_n\epsilon_0}\tilde{\psi}_y^{R_n},
\]
to ensure
\[
\psi_{n, y}(x)\geq -\frac{1}{\epsilon_0}\rho(x), \qquad\forall x\in\bbs^\dm.
\]
The functions $\psi_{n,y}$ also integrate to $0$ (see \eqref{eqn:int-tildephiy}), and therefore $\psi_{n, y}\in\mathcal{A}$ for all $n \geq 1$. Also, since the functional $\calF$ is homogeneous, we have 
\[
\mathcal{F}[\psi_{n,y}]=\mathcal{F}[\tilde{\psi}_y^{R_n}].
\]

Finally, using the above, along with \eqref{eqn:Ftildephi=Fphi} and \eqref{eqn:limF-Rinf}, we conclude that $\{\psi_{n, y}\}_{n\geq1}\subset\mathcal{A}$ satisfies 
\[
\lim_{n\to\infty}\mathcal{F}[\psi_{n, y}]= \lim_{n\to\infty}\mathcal{F}[\tilde{\psi}_y^{R_n}] = \mathcal{F}[\tilde{\psi}_y]=\mathcal{F}[\psi_y].
\]
\end{proof}

\end{document}
